\documentclass[../../../include/open-logic-section]{subfiles}

\begin{document}

\olfileid[tr]{sth}{z}{sep}
\olsection{Ayırma}

Naif Küme Oluşturma'nın yerine geçecek bir ilkeyle başlıyoruz:

\begin{axiom}[Ayırma Şeması] Her $\phi(x)$ formülü için şu bir
aksiyomdur: Her $A$ için $\Setabs{x \in A}{\phi(x)}$ kümesi vardır.
\end{axiom}

Bunun tek bir aksiyom olmadığına dikkat edin. Bu bir aksiyom
\emph{şemasıdır}. Her $\phi(x)$ formülü için bir tane olmak üzere
\emph{sonsuz sayıda} Ayırma aksiyomu vardır. Şema eşdeğer biçimde
(ve genellikle) şöyle yazılır:

\begin{defish}
``$S$'' içermeyen her $\phi(x)$ formülü için şu bir aksiyomdur:
\[
	\forall A \exists S \forall x(x \in S \liff (\phi(x) \land x \in A)).
\]
\end{defish}

\olref[sth][][]{part} başında belirtilen uzlaşıya uygun olarak,
Ayırma aksiyomlarındaki $\phi$ formülleri parametreler
içerebilir.\footnote{Bunun ne anlama geldiğine ilişkin bir açıklama için
\olref[sfr][infinite][induction]{natinductionschema} sonrasındaki
tartışmaya bakın.}

Ayırma, anlattığımız birikimli-yinelemeli küme anlayışıyla doğrudan
gerekçelendirilir. Nedenini görmek için $A$'nın bir küme olduğunu varsayın.
Öyleyse $A$, bir $S$ aşamasında oluşturulmuştur (\stageshier{} gereği).
$A$, $S$ aşamasında oluşturulduğuna göre $A$'nın bütün elemanları $S$
aşamasından önce oluşturulmuştur (\stagesacc{} gereği). Şimdi özellikle,
$A$'nın elemanı olup ayrıca $\phi$'yi sağlayan bütün kümeleri ele alın;
açıktır ki bu kümelerin tümü de $S$ aşamasından önce oluşturulmuştur.
Dolayısıyla bunlar $S$ aşamasında $\Setabs{x \in A}{\phi(x)}$ kümesi
olarak da oluşturulur (\stagesacc{} gereği).

Ayırma, Naif Küme Oluşturma'nın aksine Russell paradoksundan kaçınır.
Çünkü $\Setabs{x}{x \notin x}$ kümesinin varlığını doğrudan ileri
süremeyiz. Bunun yerine, verili bir $A$ kümesi için
$R_A = \Setabs{x \in A}{x \notin x}$ kümesinin varlığını ileri sürebiliriz.
Fakat bundan çıkan yalnızca $R_A \notin R_A$ ve $R_A \notin A$ sonuçlarıdır;
bunların hiçbiri kaygı verici değildir.

Ne var ki Ayırma'nın çarpıcı ve dolaysız bir sonucu vardır:

\begin{thm}\ollabel{thm:NoUniversalSet}
\emph{Evrensel} bir küme yoktur; yani $\Setabs{x}{x = x}$ mevcut değildir.
\end{thm}

\begin{proof}
Olmayana ergiyle, $V$'nin evrensel bir küme olduğunu varsayın. O zaman
Ayırma gereği $R = \Setabs{x \in V}{x \notin x} =
\Setabs{x}{x \notin x}$ vardır; bu da Russell paradoksuyla çelişir.
\end{proof}

Evrensel bir kümenin bulunmaması---aslında kümeler hiyerarşisinin açık
uçluluğu---birikimli-yinelemeli anlayışın arkasındaki en temel fikirlerden
biridir. Bu sonuca sezgisel olarak başka bir yoldan da ulaşabileceğimizi
görmek yararlıdır. Evrensel bir küme, kendisinin içinde !!a{element}
olmak zorundadır. Oysa birikimli-yinelemeli anlayışımıza göre her küme,
bütün !!{element}s ilk kez ortaya çıktıktan hemen sonraki ilk aşamada
hiyerarşide belirir. Bu da \emph{hiçbir} kümenin kendi kendisinin elemanı
olamayacağını gerektirir. Çünkü kendi kendisinin elemanı olan herhangi bir
kümenin, ilk ortaya çıktığı aşamadan hemen sonra ilk kez ortaya çıkması
gerekirdi; bu saçmadır. (\olref[sth][spine][rank]{defnsetrank} içinde,
bir kümenin ``rankı'' kavramını kullanarak bu açıklamayı nasıl daha kesin
hâle getireceğimizi göreceğiz. Ancak bunun için birkaç aksiyoma daha
ihtiyacımız olacak.)

Ayırma ile Kaplamsallık'ın birkaç sonucu daha şöyledir.

\begin{prop}\ollabel{prop:emptyexists}
Herhangi bir küme varsa $\emptyset$ vardır.
\end{prop}

\begin{proof}
$A$ bir kümeyse $\emptyset = \Setabs{x \in A}{x \neq x}$, Ayırma gereği
vardır.
\end{proof}

\begin{prop}
Her $A$ ve $B$ kümesi için $A \setminus B$ vardır.
\end{prop}

\begin{proof}
$A \setminus B = \Setabs{x \in A}{x \notin B}$, Ayırma gereği vardır.
\end{proof}

Neredeyse gelişigüzel kesişimlerin de var olduğu ortaya çıkar:

\begin{prop}\ollabel{prop:intersectionsexist}
$A \neq \emptyset$ ise
$\bigcap A = \Setabs{x}{(\forall y \in A)x \in y}$ vardır.
\end{prop}

\begin{proof}
$A \neq \emptyset$ olsun; öyleyse bir $c \in A$ vardır. Bu durumda
$\bigcap A = \Setabs{x}{(\forall y \in A)x \in y} =
\Setabs{x \in c}{(\forall y \in A)x \in y}$ olur ve bu küme Ayırma
gereği vardır.
\end{proof}

Ancak $A \neq \emptyset$ koşuluna dikkat edin: $\bigcap\emptyset$,
boşuna doğru olma nedeniyle evrensel küme olurdu ve bu da
\olref{thm:NoUniversalSet} ile çelişirdi.

\end{document}
